% Report fragment. Requires amsmath, amssymb, amsthm, booktabs, hyperref.
% This fragment deliberately does not assert an unconditional literal iPRA/Con(PA) endpoint.
\section{The combined AX11--Q--H command protocol}
\label{sec:ax11qh5-proof}

The new protocol combines three independently meaningful changes.  The
sixth logical head has its two premises exchanged; SHIFT tests the cached
formula alone; and AX stores its guarded gate, then falls through into the
existing no-push MP routine.  The last change consumes a fourth AX operand.
The combined protocol therefore needs its own proof compiler.  In
particular, the last command of a refutation word must be the POP gadget
defined below.  A final plain SHIFT can discard a visible contradiction.

This section gives a mathematical proof at the command and numerical-calculus
levels, together with the exact relative arithmetic statement it supports.
The executable checker verifies symbolic identities and its binding to the
delivered grammar.  No newly checked Lean proof of the final literal-machine
consistency equivalence is asserted.

The literal source binding starts from the supplied QH5 program.  Its old
AX return, source ordinal 3613, is replaced by the nineteen-site push at
ordinals 3613 through 3631; the following full MP field-reader entry is
3632.  The other 3,713 source instructions and their resolved successors
agree under the explicit relocation.  The source has
$3714-1+19=3732$ sites.  These finite graph identities are distinct from
the logical argument below and are checked separately.

\subsection{Numerical formulas and the exchanged calculus}

Let
\[
 p(a,b)=\frac{(a+b)(a+b+1)}2+b,
 \qquad
 I(A,B)=5p(A,B),\quad M(a,b)=5p(a,b)+1,
\]
\[
 U(a,B)=5p(B,a)+2,\qquad E(a,b)=5p(a,b)+3.
\]
Write $A\to B$ for $I(A,B)$, $N(B)=I(B,1)$,
$A\wedge B=N(I(A,N(B)))$, and $\exists_a B=N(U(a,N(B)))$.
The distinguished contradiction is $1$ and zero is a theorem.
The identity $I(0,0)=0$ is retained throughout.

The six logical families are
\begin{align*}
 h_1(D,H)&=(\exists_D U(D,H))\to H,\\
 h_2(D,A,H)&=U(D,U(A,H))\to U(A,U(D,H)),\\
 h_3(D,A,H)&=M(H,D)\to(E(H,A)\to M(A,D)),\\
 h_4(D,A,H)&=E(H,D)\to(M(A,H)\to M(A,D)),\\
 h_5(D,H)&=\exists_D E(H,D),\\
 h'_6(D,A,H)&=U(D,A\to H)\to(U(D,A)\to U(D,H)).
\end{align*}
There are also the restricted propositional family
\[
 \rho_8(D,A,H,u)=((H\to u)\to D)\to((D\to H)\to(A\to H)),
 \qquad u\leq D,
\]
the diagonal family $\Gamma(B)=I(B,U(B,B))$, and the guarded gate
\begin{align*}
 V(a,b,c)&=M(a,b)\to U(c,M(a,b)),\\
 P(a,b,c)&=c\to M(c,b),\\
 S(a,b,c)&=U\bigl(b,(M(a,b)\to P(a,b,c))
                         \wedge(M(c,b)\to c)\bigr),\\
 Q(a,b,c)&=U(c,S(a,b,c)\to N(P(a,b,c))),\\
 G(a,b,c)&=N(V(a,b,c)\to Q(a,b,c)),\qquad c\ne a,b.
\end{align*}
Its rules are modus ponens and theorem generalization
$B\mapsto U(i,B)$.  Call this numerical calculus $\mathcal D_5$.
An AX whose gate guard fails emits zero in the gate position.
This is harmless because $\rho_8(0,0,0,0)=0$.

The supplied QH5 packet already uses precisely $\mathcal D_5$.
Thus the new AX11 splice introduces no further logical axiom or radix
translation.  If one compares with the unexchanged coefficient-five
calculus $\mathcal C_5$, the only changed leaf is
\[
 h_6=U(D,A)\to(U(D,A\to H)\to U(D,H)).
\]
The fixed implicational theorem
\[
 (P\to(Q\to R))\to(Q\to(P\to R))
\]
converts $h_6$ to $h'_6$; the same theorem with $P,Q$ exchanged gives
the converse.  The supplied finite restricted-implication derivation
derives this theorem using only modus ponens, restricted $\rho_8$,
reflexivity of $\leq$, and the two inequalities
$a,b\leq 5p(a,b)$.  Alternatively the accompanying exchange certificate
replays its finite $K,S,\mathrm{MP}$ derivation, with $K,S$ supplied by
that same restricted-implication derivation.  Recursive replacement of
sixth-axiom leaves preserves every conclusion.  This local proof
translation does not require a monotonicity claim about hereditary radix
recoding.

\subsection{Exact four-field AX semantics}

A store is $(t,w,d,h)$: top formula, Cantor stack, numerical offset,
and cached formula.  Put $[a;w]=p(a,w)$, with total projections
$\operatorname{hd}0=\operatorname{tl}0=0$.  Hence $[0;0]=0$;
the following calculations never require a free-list interpretation.

For $D=d+x$, $H=h+y$, and $A=c$, define
\[
\begin{array}{c|l}
j&q_j\\\hline
0&G(H,D,A)\text{, or zero if its guard fails}\\
1&\Gamma(A)\\
2&h'_6(D,A,H)\\
3&h_5(D,H)\\
4&h_4(D,A,H)\\
5&h_3(D,A,H)\\
6&h_2(D,A,H)\\
7&h_1(D,H)\\
8&\rho_8(D,A,H,d)\\
9&U(d,h).
\end{array}
\]
The exact command effects are
\[
 \mathsf A_{c,x,y,z}(t,w,d,h)
       =(0,[q_0,q_1,\ldots,q_9,t;w],z,0),
 \tag{A11H}
\]
\[
 \mathsf{MP}_x(t,w,d,h)=
 \begin{cases}(b,w,b,0),&I(t,b)=h,\\
 (0,w,b,0),&I(t,b)\ne h,
 \end{cases}
 \qquad b=d+x,
 \tag{MPH}
\]
and
\[
 \mathsf S(t,w,d,h)=
 \begin{cases}\mathsf{halt},&h=1,\\
 (h,\operatorname{tl}w,d,\operatorname{hd}w),&h\ne1.
 \end{cases}
 \tag{SH}
\]

The added final AX push leaves $t=d=h=0$ before entering the full MP
field reader.  That reader consumes $z$ from the continuation.  If $z=0$,
MP succeeds with conclusion zero.  If $z>0$, the comparison
$I(0,z)=0$ fails, and MP again returns zero.  Both cases preserve the
eleven-entry stack, leave $d=z$, and clear $h$.  This proves (A11H) for
every fourth field, including malformed attempts to regard it as part of
the next command.  The compiler chooses $z=0$; the soundness argument
must cover all $z$.

The checker expands the actual grammar in a free constructor algebra.
It checks the entire ordered stack, the exchanged sixth head, the gate
before the guard, both guarded outputs, the final push, the coefficient-five
multiplier node, the no-push MP controller, the cache-only SHIFT controller,
and the full MP operand-reader shape.  This is stronger than checking a
finite collection of numerical parameter values.  It remains a finite
symbolic certificate interpreted by the accompanying mathematical lemmas.

\subsection{Gadgets, including all exceptional cases}

Let $A_0=\mathsf A_{0,0,0,0}$, and execute words left to right.
Every $q_0,\ldots,q_8$ has residue zero modulo five, and
$q_9=U(d,h)$ has residue two and is at least two.  Therefore none of
these generated heads equals one.  Direct iteration of (SH) gives
\begin{align*}
 \mathsf A_{c,x,y,0};\mathsf S^{10}(t,w,d,h)
     &=(q_8,[t;w],0,q_9),\\
 \mathsf A_{c,x,y,0};\mathsf S^{11}(t,w,d,h)
     &=(q_9,w,0,t),
\end{align*}
without halting, for arbitrary $t,w,d,h$.

Define
\[
 \mathsf{Push}=A_0;\mathsf S^{10};\mathsf{MP}_0,
 \qquad \mathsf{Pop}=A_0;\mathsf S^{12},
 \qquad \mathsf{Discard}=\mathsf{MP}_0,
 \qquad \mathsf{Offset}_i=\mathsf{MP}_i.
\]
Their identities are
\begin{align*}
 \mathsf{Push}(t,w,d,h)&=(0,[t;w],0,0),\\
 \mathsf{Pop}(t,[a;w],d,h)&=
  \begin{cases}\mathsf{halt},&t=1,\\(t,w,0,a),&t\ne1,\end{cases}\\
 \mathsf{Discard}(t,w,0,h)&=(0,w,0,0),\\
 \mathsf{Offset}_i(0,w,0,0)&=(0,w,i,0).
\end{align*}
For Push, the final comparison fails because its major candidate is an
implication code whereas its cache is a universal code.  Discard has
the same conclusion zero in both MP branches.  Pop's last SHIFT tests
exactly the saved top $t$; in the halting case it does so before reading
the old underlying stack.  These equations include empty stacks and
all collapsed zero pairs.

For $0\leq j\leq8$, the selector
\[
 \mathsf{Discard};\mathsf A_{A,D,H,0};
 \mathsf S^{j+2};\mathsf{Pop}^{9-j}
 \tag{SEL11H}
\]
sends every $(t,w,0,0)$ to $(q_j,w,0,0)$.
Each Pop sees the selected $q_j\ne1$, so all the indicated Pops
complete.  After the initial Shifts the remaining generated suffix
and the saved zero have exactly $9-j$ stack cells.  For
$\rho_8(D,A,H,u)$ with $u\leq D$, use instead
\[
 \mathsf{Discard};\mathsf{Offset}_u;
 \mathsf A_{A,D-u,H,0};\mathsf S^{10};\mathsf{Pop}.
\]
This emits an actual restricted leaf with its required inequality;
the numerical offset is never silently set to a formula parameter
whose ordering has changed.

Generalization uses
\[
 \mathsf{Gen}_i=\mathsf{Push};\mathsf{Offset}_i;\mathsf S;
                \mathsf A_{i,0,0,0};\mathsf S^{11}.
\]
It sends $(a,w,0,0)$ to $(U(i,a),w,0,0)$ for all $a,w,i$.
The intervening SHIFT has cache zero.  The old formula $a$ becomes an
AX parameter before another SHIFT tests it; in particular this identity
also holds when $a=1$.

Finally, on $(0,[a,I(a,b),0;w],0,0)$ the word
\[
 \mathsf S;\mathsf S;\mathsf{MP}_b;\mathsf{Pop}
 \tag{MP11H}
\]
halts if $a=1$ or $b=1$, and otherwise returns $(b,w,0,0)$.
The first two Shifts expose the minor and major premises.  MP leaves
the zero padding below them, and Pop removes that padding while restoring
clean caches.  The padding is essential because the MP routine itself
no longer pushes a dummy zero.  When all displayed numerical values are
zero, total projections give the same result despite $p(0,0)=0$.

\subsection{A compiler valid over every carried store}

Compile a finite $\mathcal D_5$ derivation $e$ to a word $W_e$.
Use the selectors for axiom leaves.  For a generalization append
$\mathsf{Gen}_i$.  At modus ponens, with major derivation $e_1$ of
$I(a,b)$ and minor derivation $e_0$ of $a$, use
\[
 W_e=\mathsf{Discard};\mathsf{Push};W_{e_1};\mathsf{Push};
 W_{e_0};\mathsf{Push};\mathsf S;\mathsf S;
 \mathsf{MP}_b;\mathsf{Pop}.
 \tag{COMP11H}
\]
Induction on the derivation proves the exact invariant:
from every $(t,w,0,0)$, this word either halts or returns
$(\operatorname{concl}(e),w,0,0)$.  In the MP case the initial
Discard and Push reserve zero; the two recursively compiled words
preserve the entire protected old stack; (MP11H) then applies.
The early-halting cases satisfy the invariant without any consistency
assumption.

If $e$ concludes one, the word
\[
                 A_0;W_e;\mathsf{Pop}
 \tag{HALT11H}
\]
halts from every store $(t,w,d,h)$.  The first AX clears both caches
and creates a legitimate new underlying stack.  The invariant applies
over that entire stack.  If no earlier command halts, the final Pop
tests the returned one before reading below the protected part.

This last command is a substantive change from the sum-guard AX11
compiler.  Here
\[
 \mathsf S(1,0,0,0)=(0,0,0,0),
\]
so replacing the final Pop by a plain SHIFT destroys the argument.
The checker tests this failure explicitly rather than relying on the
name of an inherited gadget.

\subsection{Why the command argument works in weak arithmetic}

The preceding arbitrary-stack invariant is convenient mathematical
notation, but one must not apply $\Sigma_1$ induction to a formula of
the form ``there is a word that works over every stack.''  The following
tracked execution supplies a $\Sigma_1$ invariant.

Run a word from the zero store with a height counter $\delta=0$.
AX increases $\delta$ by eleven, MP leaves it unchanged, and a
nonhalting SHIFT decreases it by one.  A nonhalting SHIFT at height
zero raises an underflow flag.  A halting SHIFT tests the cache before
reading the stack, so it needs no positive height.  Pair the explicit
finite list above the protected base with its height; keeping the
height separately accommodates collapsed numerical zero entries.

For fixed word, trace, and base stack $w$, induction on its coded trace
shows that a run without underflow is independent of $w$ below the
tracked segment.  The top, offset and cache agree, and the stack is
the tracked finite segment folded over $w$.  This is a direct invariant
on coded finite sequences.  All its checks are primitive recursive.
The universal quantifier over $w$ is introduced after this proof;
it is absent from the induction predicate.

Let $\operatorname{Good}(r)$ assert that there are a word and a tracked
trace which, from the zero store, avoid underflow and either halt or
finish at $(r,0,0,0)$ at height zero.  It is a $\Sigma_1$ predicate in
arithmetic.  The gadgets give the following height ledger:
\[
\begin{array}{l|c}
\text{gadget}&\text{net height on a nonhalting return}\\\hline
\mathsf{Push}&+1\\
\mathsf{Pop}&-1\\
\text{selector }j&11-(j+2)-(9-j)=0\\
\mathsf{Gen}_i&1-1+11-11=0\\
\text{MP compiler frame}&3-2-1=0.
\end{array}
\]
Every prefix of a compiled nonhalting word has nonnegative tracked
height.  Each compiled word begins with Discard, so its initial top
is irrelevant whenever the initial caches are clean.  The frame lemma
therefore supplies exactly the arbitrary-store property needed for
composition.  The fixed selectors prove $\operatorname{Good}$ for
every leaf; the displayed words prove its closure under generalization
and MP.  $\Sigma_1$ induction along an internally coded derivation, retaining
an existentially coded list of the preceding word-and-trace witnesses,
proves $\operatorname{Good}(r)$ for every internally provable $r$,
including nonstandard proof codes in a model of $\mathrm I\Sigma_1$.
Appending the final Pop to a successful proof of one halts before
underflow.

This is the same tracked-executor method used by the repository's
\texttt{pa194-model-fairness} development, with the eleven-push AX
contract and the new terminal Pop substituted explicitly.  The present
result is a paper proof; the existing ten-entry Lean theorem is not
being asserted to have checked this new interface.

\subsection{Four-field scheduling and soundness}

There are two different integer conventions at the fetch boundary.
They can be read directly from the literal source instructions 25--47,
which the checker binds independently.  Register 0 is the next fresh
candidate $N$, register 8 is the continuation $L$, and register 3
receives the command tag.  Registers 3, 6 and 7 are initially zero,
as required at a clean command-fetch boundary.  At source site 25 the
seed to be unpaired is
\[
 \operatorname{seed}(N,L)=
 \begin{cases}N,&L=0,\\L-1,&L>0,\end{cases}
 \qquad
 N'=\begin{cases}N+1,&L=0,\\N,&L>0.\end{cases}
 \tag{SEED11H}
\]
Indeed, positive $L$ is decremented once and execution jumps directly
to site 34.  For $L=0$, sites 27--30 copy $N$ into $L$ and scratch 6;
sites 31--33 restore $N$ with one additional unit and clear scratch 6.
Sites 34--47 then unpair the seed, putting its first component in
register 3 and its second in register 8.  At the head of their main
loop the invariant is
$\operatorname{seed}=p(R_6,R_7)+R_8$.  One outer iteration replaces
$(a,b)$ by $(a-1,b+1)$ when $a>0$, or by $(b+1,0)$ when $a=0$;
in either case its Cantor pair increases by one while $R_8$ decreases
by one.  The two concluding transfers clear the scratches and return
the two projections.  This proves (SEED11H) for arbitrary natural
inputs, preserves the other registers, and also supplies its internal
finite arithmetic induction proof.  The initial clean frame has
$N=2,L=0$.

The command continuation must retain all four fields.  Define
$e(\epsilon)=0$ and $e(c::w)=1+b_c(e(w))$, where
\begin{align*}
 b_{\mathsf S}(v)&=p(0,v),\\
 b_{\mathsf{MP}_x}(v)&=p(1,p(x,v)),\\
 b_{\mathsf A_{c,x,y,z}}(v)&=p(c+2,p(y,p(x,p(z,v)))).
\end{align*}
Thus a \emph{fresh candidate} for $c::w$ has value $b_c(e(w))$,
without an additional one, whereas a \emph{positive continuation}
for the same word has value $e(c::w)=1+b_c(e(w))$.
Equation (SEED11H) makes both enter the same decoder with payload
$b_c(e(w))$.  This is why the candidate family below has no outer
``$+1$''.  The fields after the command tag are read in the order $y,x,z$.
Every inverse equation follows from total Cantor projections, including
zero aliases.  The checker also decodes the old three-field encoding
and confirms that it wrongly consumes the head of the next continuation
as $z$.

After decoding a positive pointer $1+b_c(v)$, its continuation is a
component of $b_c(v)$ and is strictly smaller than the old pointer.
Consequently each candidate executes finitely many commands unless it
halts.  Fresh candidate numbers start at two and increase.  If
$v=e(A_0;W_e;\mathsf{Pop})$, then prefixing by
$\mathsf A_{k,0,0,0}$ gives candidate codes
\[
       p(k+2,p(0,p(0,p(0,v)))),\qquad k\geq0.
\]
This is an increasing unbounded family, with every member at least
three.  Every such candidate halts because its suffix works from
every store, regardless of what its prefix and earlier candidates
left behind.  Bounded induction on the earlier candidate numbers
proves that the scheduler reaches one or halts earlier.  The decrease
of the continuation gives a concrete command-count bound: a fresh
candidate $n$ executes at most $n+1$ commands, since its first remainder
is at most $n$ and every positive remainder subsequently decreases.
Thus candidates $2,\ldots,r$ together execute at most
\[
 \sum_{n=2}^{r}(n+1)=\frac{r^2+3r-4}{2}\leq r^2
 \qquad(r\geq2)
\]
commands.  Take the $k=0$ member of the displayed family as $r$.
In the command-step scheduler, bounded minimization through $r^2$
steps therefore finds a halt whenever the input is a valid refutation.
The candidate code, total command functions, bounded execution and
bounded minimization are all primitive recursive.  This supplies an
explicit witness transformation; the bound counts commands and does
not assert a bound of $r^2$ on literal Turing-machine transitions.
No unbounded search for a successful candidate is used.

For the converse, annotate the top, cache and every canonical
Cantor-stack entry with a $\mathcal D_5$ derivation.  The zero store
has such an annotation.  At an AX, each $q_j$ is an axiom or, for
$q_9$, the generalization of the annotated old cache; the saved old
top retains its annotation.  The restricted leaf has
$d\leq d+x$.  The arbitrary fourth operand carries no formula
obligation and the two new visible zero entries have their fixed
proofs.  A successful MP composes its annotated premises; a failed
MP produces zero.  SHIFT moves annotations and halts only when
the annotated cache is exactly one.  For an internally finite
trace, the existence of its coded annotation history is a
$\Sigma_1$ induction invariant.  This yields an internal
$\mathcal D_5$ refutation from every scheduler halt.

Thus, for the explicit four-field scheduler and the displayed
primitive-recursive numerical proof checker, the preceding paper
argument establishes
\[
 \mathrm I\Sigma_1\ \vdash\
  H_{\mathrm{AX11QH5,sch}}
    \longleftrightarrow
  \exists e\,\operatorname{Prf}_{\mathcal D_5}(e,1).
 \tag{REL-SCH}
\]
Both witness conversions are primitive recursive: the forward one
maintains proof annotations, and the reverse one constructs the word,
its candidate code, a bounded scheduler trace, and its first halting
index by bounded minimization.  The proof specifies its weak-arithmetic
induction predicates; it does not infer (REL-SCH) merely from a theorem
over the standard natural numbers.

\subsection{What remains conditional at the literal endpoint}

Let $T$ be the delivered binary table and let $H_T$ be its arithmetized
blank-tape halting sentence.  Two distinct statements complete the
requested iPRA result:
\begin{align*}
 \mathrm I\Sigma_1&\vdash H_T\leftrightarrow
                        H_{\mathrm{AX11QH5,sch}},\tag{PHYS}\\
 \mathrm I\Sigma_1&\vdash
  \exists e\,\operatorname{Prf}_{\mathcal D_5}(e,1)
       \leftrightarrow \operatorname{Ref}_{\mathrm{PA}}.\tag{CAL}
\end{align*}
The exact finite startup, local folds and universal symbolic register
contracts are evidence for (PHYS); establishing its internal coded
simulation requires the uniform frame, source-controller and positive
runtime inductions.  A conditional proof should keep those obligations
visible.  The inherited manuscripts state PA-provable calibration;
their PA proofs do not automatically establish (CAL).

Given (PHYS) and (CAL), composition with (REL-SCH) gives the all-model
calibration required by the repository's
\texttt{ConPAUnconditional.machine\_endpoint}.  The established
conservation and negative-translation pipeline then yields
\[
 \mathrm{iPRA}\ \vdash\
 \operatorname{Con}(\mathrm{PA})\leftrightarrow\neg H_T.
\]
Without those premises this last formula remains conditional.  In
particular, exchanging an axiom head, changing the command protocol,
or verifying every ordinary finite run cannot supply an unproved
uniform PA-refutation compiler.

At the relative command level, (REL-SCH) and the analogous theorem
for the original QH5 scheduler show that the AX11 splice preserves
the existence of a numerical-calculus refutation.  That is the useful
unconditional contribution of the present protocol proof.  The same
gadget proof works for every coefficient $k\geq4$: it uses only
the zero alias, the disjoint implication/universal residues, and
$U(a,B)\geq2$.  The delivered checker is bound to the actual
coefficient-five source, and the coefficient-four case is not
being advertised as a second delivered machine.
